\documentclass[11pt]{amsart}

\usepackage{amsmath,amssymb,amsthm,mathtools}
\usepackage{enumitem}
\usepackage[hidelinks]{hyperref}
\usepackage{microtype}

% PeAIce claim-discipline macros (input from manuscript preambles)
% Status tags for theorem environments and inline use

\newcommand{\PeAIceFormal}{\textsc{Formal}}
\newcommand{\PeAIceProposed}{\textsc{Proposed}}
\newcommand{\PeAIceOpen}{\textsc{Open}}
\newcommand{\PeAIceKnown}{\textsc{Known}}
\newcommand{\PeAIceNumerics}{\textsc{Numerics}}
\newcommand{\PeAIceClosedNeg}{\textsc{Closed-Negative}}
\newcommand{\PeAIceClosedPos}{\textsc{Closed-Positive}}
\newcommand{\PeAIceAnalogy}{\textsc{Structural Analogy}}

\newcommand{\statusbox}[1]{%
  \par\noindent\fbox{\parbox{0.95\linewidth}{\small\textbf{Status.} #1}}\par\medskip
}

\newcommand{\firewallbox}[1]{%
  \par\noindent\fbox{\parbox{0.95\linewidth}{\small\textbf{Firewall.} #1}}\par\medskip
}

\newcommand{\PeAIceOpenTargets}{%
  The {R}iemann hypothesis and the {C}oleman conjecture remain \PeAIceOpen.
  No claim in this note constitutes a proof of either statement.
}

\newcommand{\ZetaFirewall}{%
  The value $\zeta(0)=-\tfrac12$ is \emph{not} asserted as a term in the
  {K}akeya proof literature. Any regularization at $s=0$ is second-stage and
  conditional on a residual operator or counting object constructed after
  residual control.
}

\newcommand{\KNSFirewall}{%
  Typed incidence objects do not locate zeros of $\zeta$.
  Overlap measures are not placement measures on the critical line.
}


\theoremstyle{plain}
\newtheorem{theorem}{Theorem}[section]
\newtheorem{lemma}[theorem]{Lemma}
\newtheorem{proposition}[theorem]{Proposition}

\theoremstyle{definition}
\newtheorem{definition}[theorem]{Definition}
\newtheorem{remark}[theorem]{Remark}

\title[KakeyaNeedleSet(Light(Basic))]{KakeyaNeedleSet(Light(Basic)):
a typed incidence object and two-layer decomposition}
\author{Manuel Coleman}
\address{Love Labs LCA / PeAIce Research Program}
\email{coleman.manuel23@gmail.com}
\date{\today}

\begin{document}

\begin{abstract}
We introduce a typed incidence object
$\mathrm{KakeyaNeedleSet}(\mathrm{Light}(\mathrm{Basic}))$ (abbreviated
$\mathrm{KNS}(\mathrm{LB})$): a single throughput center with a full
directional fan of unit segments in dimensions $d\in\{2,3\}$, together with
its $\delta$-tube discretization. We record elementary fan geometry and a
two-layer separation between \emph{seen} overlap statistics (multiplicity,
clustering) and \emph{unseen} center-action / placement declarations.
Overlap functionals do not determine lawful placement and do not locate zeros
of the Riemann zeta function.
\PeAIceOpenTargets\
\KNSFirewall
\end{abstract}

\subjclass[2020]{42B25, 28A75}
\keywords{Kakeya sets, tubes, multiplicity, incidence geometry, typed objects}

\maketitle

\section{Introduction}

\statusbox{Program note ARX-002. Definitions and fan lemmas are \PeAIceFormal
in stated scope. Any bridge to zeros is \PeAIceOpen. Compute receipts from
program probes are \PeAIceNumerics and not theorems of $\zeta$.}

Modern Kakeya proofs import a precise vocabulary of tubes, shadings, and
multiplicity \cite{WangZahl2025Kakeya,GuthWangZahl2026Streamlined}. This note
isolates a \emph{measure-maximal} configuration---the full fan through a
center---as a typed object useful for separating incidence observables from
placement declarations in a larger research program. It is not a proof of the
Riemann hypothesis.

\section{Light(Basic) and KakeyaNeedleSet}

\begin{definition}[Light(Basic)]
\label{def:light}
A \emph{Light(Basic)} object in dimension $d\in\{2,3\}$ is a pair
$L=(C,\Phi_C)$ with center $C\in\mathbb{R}^d$ and intensity $\Phi_C>0$.
The center is a declared throughput point in carrier space; it is not
identified with the complex origin, nor with the value $\zeta(0)$.
\end{definition}

\begin{definition}[KakeyaNeedleSet]
\label{def:kns}
For $L=(C,\Phi_C)$ and direction sphere $S^{d-1}$ (directions up to sign),
\begin{equation}
  \mathrm{KakeyaNeedleSet}(L)
  \;:=\;
  \{\,\sigma_\omega : \omega\in S^{d-1}\,\},
  \qquad
  \sigma_\omega = \{\,C+t\omega : |t|\le \tfrac12\,\}.
\end{equation}
We write $\mathrm{KNS}(\mathrm{LB})$ for
$\mathrm{KakeyaNeedleSet}(\mathrm{Light}(\mathrm{Basic}))$.
\end{definition}

\paragraph{$\delta$-discretization.}
Let $\Omega_\delta\subset S^{d-1}$ be a maximal $\delta$-separated set, so
$|\Omega_\delta|\asymp \delta^{-(d-1)}$, and let $\mathcal{T}^C_\delta$ be the
corresponding family of $\delta$-tubes about the segments $\sigma_\omega$ with
full shading $Y(T)=T$. Multiplicity $\mu$, union $U$, and clustering quantities
are imported from the Kakeya tube calculus
\cite{GuthWangZahl2026Streamlined}.

\begin{lemma}[Fan geometry]
\label{lem:fan}
Let $\mathcal{T}^C_\delta$ be as above with full shading.
\begin{enumerate}[label=(\alph*)]
  \item $U(\mathcal{T}^C_\delta)\supset B(C,1/2)$ and $|U|\asymp 1$.
  \item Pointwise multiplicity satisfies
  $\mu(x)\asymp (\delta+|x-C|)^{-(d-1)}$ for $|x-C|\le 1/2$, hence
  $\sup_x \mu(x)\asymp \delta^{-(d-1)}$ on the core $B(C,\delta)$.
  \item Averaged multiplicity is $\asymp 1$; the configuration is the sticky,
  positive-measure extreme of the Kakeya family at this scale.
\end{enumerate}
\end{lemma}

\begin{proof}[Proof sketch]
(a) Every ray of length $1/2$ from $C$ lies in some $\sigma_\omega$ up to
$\delta$; total volume scales as
$\delta^{-(d-1)}\cdot\delta^{d-1}\asymp 1$.
(b) $T_\omega\ni x$ forces direction within angular width $\asymp\delta/|x-C|$;
counting $\delta$-caps yields the displayed decay.
(c) Average the radial profile against the natural measure on annuli.
\end{proof}

\section{Two-layer separation}

\begin{proposition}[Overlap does not determine placement]
\label{prop:sep}
Incidence data $(\mu,\Delta_{\max},U)$ generated by a tube configuration do not
determine a Hilbert-space placement functional $\ell_{\mathrm{off}}$ relative
to a declared register $A$. In particular, high core multiplicity (glare) is
compatible with total misplacement, and high average multiplicity without a
common point is compatible with aligned placement.
\end{proposition}

\begin{proof}[Proof outline]
(i) Incidence objects are geometric; placement requires a declared encoding
space and register. (ii) Fix incidence data and vary the register so that a
fixed encoding has arbitrary off-register mass. (iii) Fix encoding and rotate
the register. Either freedom kills any map from incidence alone to
$\ell_{\mathrm{off}}$. Quantitative countermodels (full fan vs.\ Besicovitch
families) appear in the program source
\texttt{PEAICE-CLAUDEV6-KNS-LB-PAPER-001}; they will be expanded in a later
draft of this note.
\end{proof}

\begin{remark}[Non-inference to zeros]
Any implication from multiplicity to zero placement on the critical line
requires an explicit map. No such map is constructed here. Bounded
square-difference coupling corridors in the companion operator program are
recorded as \PeAIceClosedNeg\ for determinant-level identification and do not
supply the missing rung.
\end{remark}

\section{What remains open}

\begin{enumerate}[leftmargin=*]
  \item Theorem-level lift of $\mathrm{KNS}(\mathrm{LB})$ beyond typed
  incidence geometry.
  \item An explicit prime-carrying operator from placement data to zero
  location (the missing rung of the companion program).
  \item Independent measurement obligations in the compute layer (program
  item CP-004) --- \PeAIceOpen, not a substitute for proofs.
\end{enumerate}

\PeAIceOpenTargets\
\KNSFirewall

\section*{Acknowledgments}
Research engineering assistance was used in preparing program materials.
All mathematical claims remain the author's responsibility.

\bibliographystyle{amsplain}
\bibliography{../../bibliography/peaice-arxiv}

\end{document}
